\section{Discussion}
\label{header:discussion}

\paragraph{}The empirical results of Section~\ref{header:results} substantiate the central hypothesis stated in Section~\ref{sect:meth4}: that a $(t, n)$ threshold-signature protocol grounded in BLS12-381 can deliver the security guarantees of multi-party authorisation without sacrificing the latency profile that users have come to expect from QR-based mobile payments. This section interprets the four streams of evidence, i.e., functional correctness, performance, adversarial resilience, and user acceptance, and places them in dialogue with the literature reviewed in Section~\ref{header:rrl}. It is organised around five axes of interpretation, each of which corresponds to a thread that recurs across the reviewed gaps in the literature.

\subsection{Cryptographic Correctness vs.\ Operational Security}
\label{sect:disc_crypto}

\paragraph{}A useful way to read the functional and adversarial results in tandem is to separate \textit{cryptographic} guarantees, which are unconditional and follow from the security of BLS and the hardness of the elliptic-curve discrete-log problem, from \textit{operational} guarantees, which depend on how the protocol is wired into the surrounding software and on the trust assumptions placed on the proxy and the persistence layer. The 35-test functional suite (Section~\ref{sect:results_func}) and the offline cryptographic attack group A1--A7 (Section~\ref{sect:results_mitm}) jointly establish the cryptographic guarantees: every legitimate 3-of-5 aggregate verifies, no 2-of-5 partial aggregate verifies, no payload mutation goes undetected, and no off-curve point is admitted by the verifier. These results are deterministic --- they hold for every input the test suite supplies, and would continue to hold under any extension of that suite that respects the protocol's stated message format.

\paragraph{}The online attack group B1--B6 instead exercises operational guarantees. The fact that the verifier rejects a forged \texttt{amount} field (B2, B3) is not a property of BLS as a primitive; it is a property of the patched \texttt{socketHandler.js}, which now reads the authoritative amount from the persisted \texttt{Transaction} document rather than from the client envelope. The original implementation, in which the client-supplied \texttt{amount} was used as the operand of the deduction, would have passed every functional test and every offline cryptographic test, and would still have admitted a non-trivial economic attack against any deployment that lacked transport-layer encryption. This separation, between cryptographic and operational guarantees, is the principal pedagogical takeaway of the present study: a threshold-signature protocol is necessary but not sufficient for an end-to-end secure payment system, because the integrity of the bound transaction \textit{message} depends on the surrounding application logic preserving the field whose value is being authorised. Tests B2 and B3 now provide regression coverage for that property.

\paragraph{}A complementary observation concerns the proxy. By construction, the proxy in this study learns nothing of the partial signatures it relays beyond their existence and timing, since the inter-client traffic is encrypted with AES-256-GCM keyed by an ephemeral Diffie--Hellman exchange (Section~\ref{sect:meth2_b}, item~h). The proxy is therefore consistent with the ``honest-but-curious'' adversary model used in much of the proxy re-encryption literature~\cite{chen2011, ruffing2014}: it can deny service and observe metadata, but cannot read or forge cryptographic content. The audit metadata it persists (session identifiers, signature counts, timestamps) is exactly the residue that the threat model permits.

\subsection{Performance Margins and the 200~ms Acceptance Criterion}
\label{sect:disc_perf}

\paragraph{}The latency benchmark substantially exceeds the $< 200$~ms acceptance criterion. The mean verification latency of $4.83$~ms and the 99th-percentile worst case of $8$~ms leave a margin of more than an order of magnitude against the criterion, even on the localhost build. Two consequences of this margin deserve interpretation.

\paragraph{}First, the protocol leaves substantial headroom for additional cryptographic work without breaching the user-perceptible latency boundary. The verification step is cheap because it is dominated by three near-constant-cost operations: a single \texttt{findOne} lookup on the \texttt{Transaction} collection, a single \texttt{save} that updates the terminal status, and a single broadcast to the room. If, in a future iteration, the verifier were extended to perform additional cryptographic checks at scan time --- for instance, a freshness proof bound to the merchant's session, a verifiable random function-based receipt, or a per-scan zero-knowledge attestation --- the cumulative cost of those additions could grow by tens of milliseconds before the 200~ms criterion would come under threat. The protocol's performance budget is therefore not pinched; it is comfortably absorbing.

\paragraph{}Second, the QR-generation cost (mean $24.73$~ms, p99 $32$~ms) is dominated by the cryptographic primitives themselves rather than by the persistence layer. This is empirically inferred from the relative sizes of the verification and generation latencies: verification involves the same number of MongoDB writes as generation, yet costs roughly five times less, because it does not perform a partial signing, an aggregation, or a self-verification pairing. The implication is that the principal lever for further latency reduction on the generation path is the cryptographic library, not the database. The \texttt{@noble/curves} library used in this study is a pure-JavaScript implementation of BLS12-381; a native (e.g., \texttt{node-gyp}-bound) implementation, or a WebAssembly port, would be expected to reduce QR-generation latency by a meaningful factor, although given the existing margin against the criterion the engineering case for that work is weaker than it would otherwise be.

\paragraph{}A third observation, anticipating Section~\ref{sect:disc_qr}, concerns which latency the user actually perceives. The total end-to-end latency observed in the benchmark (mean $29.57$~ms, p99 $36$~ms) accounts only for the \textit{computational and network} components of the round trip. The UAT, however, showed that the latency users worry about is not this backend round trip but the dynamic QR's two-minute expiry window: in Cycle~1, four of five participants asked for a clearer indication of remaining QR validity (Section~\ref{sect:results_uat1}). The backend may resolve a transaction in tens of milliseconds, yet the user's felt experience is governed by a multi-second budget --- the expiry countdown --- and by whatever fraction of it the camera consumes acquiring focus. The dominant latency in the system, from the user's standpoint, is therefore not cryptographic but temporal-perceptual, which is precisely why the Phase-4 refinement targeted the expiry timer and the QR payload rather than the cryptographic library.

\subsection{The Security--Scannability Trade-off Between QR Encodings}
\label{sect:disc_qr}

\paragraph{}The iterative UAT pipeline produced two distinct findings about the dynamic QR code, and it is worth separating them, because the Cycle-1 finding is strong while the Cycle-2 finding is deliberately reported as modest. The Cycle-1 finding (Section~\ref{sect:results_uat1}) was that the dominant interaction friction was the QR-expiry window: four of five participants asked, unprompted, for a clearer indication of remaining validity. This was unambiguous and motivated the Phase-4 visible-timer refinement, whose success is visible in Cycle~2, where no participant reported ``high pressure'' on the expiry item and none raised the timer in open-ended feedback. The Cycle-2 finding concerns the QR \textit{encoding} itself, and is the subject of this subsection.

\paragraph{}The Big QR and Small QR encodings (contrasted in Figure~\ref{fig:qr-comparison} of Section~\ref{sect:results_refine}) sit at opposite ends of a security--scannability axis. The Big QR places the entire authorisation envelope --- $\{PK, g, F_{\text{sig}}, M\}$ --- inside the QR matrix. This makes the variant \textit{self-contained}: a verifier holding only the QR can validate the authorisation without contacting any third party, and there is no separate retrieval channel for an adversary to compromise. The cost is data density. BLS12-381 group elements are individually compact (a $G_1$ point in compressed form is 48 bytes), but the joint public key, generator, and signature together with the transaction message bring the serialised payload well beyond what a typical payment QR --- which usually carries only a URL or a short token --- would hold, forcing a higher QR version with a denser module grid~\cite{worldbank2021ch5, corpuz2025}. The denser the grid, the more precisely a camera must focus to resolve adjacent modules.

\paragraph{}The Small QR, by contrast, embeds only a short reference token; the verifier retrieves the full envelope out of band from the proxy at scan time. This lowers the module density substantially, at the cost of self-containment: a compromised retrieval channel could in principle substitute a forged envelope. In practice that channel is the same TLS-protected proxy the rest of the protocol already trusts, so the additional surface is small --- but it is real, and it shifts a portion of the trust budget from cryptographic content to transport.

\paragraph{}The UAT evidence for this trade-off is consistent in direction but modest in magnitude, and the discussion is careful not to overstate it. Item~4.2 (Section~\ref{sect:results_uat_qr}) recorded one Big QR participant out of five reporting a 2--3 second focus delay, against zero of five for the Small QR; every other participant in both groups reported immediate recognition, and no participant reported a hard scan failure. With five participants per group this rests on a single observation, and it is reported as a modest confirmation of a well-documented optical effect rather than as a strong empirical contrast. What the UAT does establish firmly is the weaker but still useful claim that the Big QR did not \textit{break} scannability: even the in-band encoding was recognised on the first attempt by four of five participants and never failed outright.

\paragraph{}Because the trade-off is genuine but the populations are small, the protocol design does not force a global choice between the two encodings. The threshold $t$ is already a per-transaction policy decision keyed by transaction value (Section~\ref{sect:meth2_b}, item~d), and the same Policy Engine can select the QR encoding on the same basis. This is the recommendation developed in Section~\ref{sect:rec_adaptive}: retain both encodings and let transaction value decide between the self-contained Big QR and the faster-scanning Small QR.

\subsection{Perceived Security and the Multi-Party Mental Model}
\label{sect:disc_perceived}

\paragraph{}A consistent finding in the perceived-security literature is that users underweight cryptographic mechanisms that they do not directly observe and overweight procedural mechanisms that they do~\cite{krombholz2014}. The UAT in the present study is consistent with that pattern: every participant rated the multi-party approval workflow as ``significantly safer'' or ``somewhat safer'' than a single-user payment, but no participant referred to BLS, threshold signatures, or pairing-based verification when asked to articulate their reasoning. The mechanism that participants saw, which is a pop-up on each approver's phone, a visible quorum, and a refusal to proceed without it, dominated the mechanism they did not see (the cryptographic enforcement that made the visible mechanism load-bearing).

\paragraph{}This has two consequences for the design. First, the visible enforcement of the quorum is a security-critical UI element in its own right; degrading it (for instance, by allowing approvals to be silently auto-confirmed under any circumstance, or by hiding the approver list from the shopper) would be expected to materially reduce perceived security even if the underlying protocol were unchanged. Second, the open-ended feedback's request for stricter credential confirmation at the point of approval (Section~\ref{sect:results_uat_open}) is consistent with this pattern: participants want a visible per-approval check that ratifies the cryptographic step they cannot see. Recommendation~\ref{sect:rec_mfa} is therefore as much a perceived-security recommendation as it is a credential-binding one.

\paragraph{}It is worth noting, however, the limit of this finding. None of the participants in this UAT had cryptographic training; the perceived-security ratings should be read as evidence about lay users' mental models, not about the security properties of the protocol itself. The protocol's security claims rest on the mathematical and adversarial results of Sections~\ref{sect:results_func} and~\ref{sect:results_mitm}, not on user-reported confidence.

\subsection{Limitations and Threats to Validity}
\label{sect:disc_limits}

\paragraph{}Four principal limitations should be flagged for the integrity of the foregoing claims.

\paragraph{a.) Sample size and the SUS divergence.}
The UAT was conducted across two cycles with fifteen participants in total --- five in Cycle~1 and ten in Cycle~2, the latter split into two disjoint groups of five --- with no participant taking part in more than one session. This is adequate for a formative usability evaluation in the small-sample tradition~\cite{brooke1996sus, bangor2008sus} but is not sufficient for inferential claims about population-level usability or perceived security. Two consequences follow. First, the between-group SUS divergence in Cycle~2 (Group~A mean $95.0$, Group~B mean $60.8$; Section~\ref{sect:results_uat_sus}) cannot be attributed to the QR encoding: the groups are disjoint and small, the Small QR group scanned at least as fast as the Big QR group, and the Group~B mean is pulled down by two participants whose low ratings accompanied substantive critiques rather than usability failures. The SUS is therefore reported as a single $n = 10$ aggregate of $77.9$. Second, the QR-recognition contrast (one delayed scan of five Big QR trials, against zero of five Small QR) rests on a single observation and is treated throughout as a modest, direction-only signal rather than a point estimate of any underlying rate. The six-item SUS adaptation also departs from the canonical 10-item instrument, so the resulting score is not strictly interchangeable with normative SUS benchmarks; it is reported in Section~\ref{sect:results_uat_sus} with that caveat.

\paragraph{b.) Latency benchmark scope.}
The latency benchmark was conducted on localhost, eliminating network round-trip time as a variable. The production deployment introduces an additional 50--150~ms of geographic routing overhead, which the discussion in Section~\ref{sect:disc_perf} addresses qualitatively. A controlled benchmark on the production environment, or on a representative cloud-to-mobile path, would be a useful confirmation. In addition, the benchmark exercises the low-value (threshold $t = 1$) path; the high-value path ($t \geq 2$) introduces additional inter-client coordination, whose latency profile is bounded above by the slowest approver's response time and is therefore dominated by human factors rather than by cryptography.

\paragraph{c.) Threat model boundary.}
The MITM evaluation models a network-level adversary who has either intercepted plaintext (in non-TLS deployments) or broken TLS (in production deployments), and a cryptographic adversary who can submit arbitrary candidate signatures and message digests to the verifier. It does not model an endpoint-compromise adversary who has obtained one of the five private key shares from a participant device; in that scenario, the threshold property continues to provide protection (a single share is insufficient to forge), but a coalition of $t$ compromised devices would, by construction, be able to authorise transactions. Defences against the endpoint-compromise scenario, such as proactive secret sharing, hardware-backed key storage, and behavioural anomaly detection on signing patterns, are out of scope for the present study and are noted in Section~\ref{sect:rec_proactive} and Section~\ref{sect:rec_hardware} as directions for follow-up work.

\paragraph{d.) Sub-threshold transaction splitting.}
A Cycle-2 participant observed that, because transactions below ₱1{,}000 are authorised at threshold $t = 1$ (the requester alone), a malicious co-owner could drain a joint balance by issuing many sub-threshold payments, each individually below the quorum trigger (Section~\ref{sect:results_uat_open}). This is a correct observation and a genuine limitation of the present policy: the threshold is evaluated per transaction, with no cumulative or velocity term. The protocol's cryptographic and adversarial guarantees (Sections~\ref{sect:results_func} and~\ref{sect:results_mitm}) are unaffected --- every individual sub-threshold transaction is still correctly authorised and verified --- but the \textit{policy} that maps amount to threshold does not bound aggregate sub-threshold spend over a window of time. A velocity-aware policy extension is proposed as Recommendation~\ref{sect:rec_velocity}.
